\VerdierSemanticAnchor{front-title}
\VerdierPhysicalPageAnchor{0002}
\begin{titlepage}
  \centering
  {\Huge\bfseries 239\hfill ASTÉRISQUE\par}
  \vspace{1.5em}
  {\LARGE\bfseries 1996\par}
  \vfill
  {\Huge\bfseries アーベル圏の\par}
  \vspace{0.4em}
  {\Huge\bfseries 導来圏\par}
  \vspace{3em}
  {\Large\bfseries Jean-Louis VERDIER\par}
  \vfill
  {\Large\bfseries フランス数学会\par}
  \smallskip
  {\small\bfseries フランス国立科学研究センターの助成により刊行\par}
\end{titlepage}

\VerdierSemanticAnchor{front-ams-classification}
\VerdierPhysicalPageAnchor{0003}
\chapter*{AMS 主題分類}
\addcontentsline{toc}{chapter}{AMS 主題分類}

14F20, 18E30, 18E35, 18G05, 18G10, 18G15, 18G20, 18G35,
18G40, 55U15, 55U25, 55U30.

\clearpage
\VerdierSemanticAnchor{front-portrait}
\VerdierPhysicalPageAnchor{0004}
\thispagestyle{empty}
\vspace*{0.62\textheight}
\begin{center}
Jean-Louis VERDIER、1977年
\end{center}

\clearpage
\VerdierPhysicalPageAnchor{0005}
\thispagestyle{empty}
\null
\clearpage

\VerdierSemanticAnchor{front-preface}
\VerdierPhysicalPageAnchor{0006}
\chapter*{序文}
\addcontentsline{toc}{chapter}{序文}

1989年8月25日に悲劇的な死を遂げた Jean-Louis Verdier は、1967年6月14日、
パリにおいて『アーベル圏の導来圏』と題する国家博士論文の審査を受けた。
本書のテキストは、この論文のタイプ原稿の最終版である。編者が行ったのは細部
の誤りを訂正することに限られる。学部に提出されていたのは「序論」と題する
部分だけであり、本文は審査時にも未完で、その後も未完のままであった。

この仕事の成立事情はかなりよく知られている（\emph{参照} [I]）。
Grothendieck は1960年代初頭、1958年のエディンバラ国際数学者会議での講演を
受けて構想していた双対性定理の大規模な一般化に必要なホモロジー代数の枠組み
を与えるため、導来圏という着想を得た。彼は Verdier に、その形式体系の構築
を博士論文の課題として提案した。1963年、Verdier は理論の本質的内容を要約
した成果集『導来圏、いくつかの結果（状態0）』を IHÉS の謄写版として刊行した
（[SGA 4 1/2] に再録）。導来圏の概念は Grothendieck に由来するが、それに
先立って三角圏の概念を導入したのは Verdier である。安定ホモトピー論に着想
を得たそこでの公理的枠組み（\emph{参照} [P]）は、驚くほど豊かな成果をもたらす
ことになった。

Verdier が自らの仕事を包括的に記述し始めたのは、それから何年も後のことで
ある。なぜ完成させなかったのか。多くの仮説が唱えられてきたが、私は次の説明
に傾く。欠けている部分は導来関手の理論である。Verdier がそこまで書き進めた
頃には、導来圏はすでに代数幾何で広く使われ、さまざまな変形と一般化が発展
していた。Deligne 流の各点での導来可能性、非加法的導来関手（Quillen 流の
ホモトピー代数）、フィルター付き導来関手などである。その時点で、成果集が
定めた枠組みは不十分になっていた。適切な観点を見いだすには新たな基礎研究が
必要だったが、当時ほかの問題に関心を向けていた Verdier には、その余裕も、
おそらく意欲もなかったのだと私は思う。時がたつにつれ導来圏に関する論文が
現れ、その利用は広がり、形式体系は数学の別分野（たとえば線形偏微分方程式）
にも浸透した。こうして Verdier の目には、この仕事を完成させることがますます
不要に映るようになったのである……。

\clearpage
\VerdierPhysicalPageAnchor{0007}
未完ではあるものの、Verdier のこのテキストは貴重で、まさに代えがたい参考
文献である（たとえば、ここに示されたスペクトル系列の独創的で優美な提示は、
私の知る限りほかのどこにも見られない）。それがついに日の目を見たことを、
心から喜びたい。

\begin{flushright}
パリ、1996年4月27日\\
Luc Illusie
\end{flushright}

\section*{参考文献}

\begin{description}[labelwidth=6.2em,leftmargin=6.8em]
\item[{[I]}] L. Illusie, \emph{Catégories dérivées et dualité, travaux de
J.-L. Verdier}, L'Enseignement Mathématique, 36 (1990), 369--391.
（1989年10月19日、Université Paris VII 主催の Jean-Louis Verdier 追悼式
における講演。）

\item[{[P]}] D. Puppe, \emph{On the formal structure of stable homotopy
theory}, Colloquium on Algebraic Topology, Aarhus Universitet (1962),
65--71.

\item[{\mbox{[SGA 4 1/2]}}] ``Cohomologie étale'', Séminaire de Géométrie
Algébrique du Bois-Marie, SGA 4 1/2, P. Deligne, Lecture Notes in
Mathematics 569, Springer-Verlag (1977), 262--311.
\end{description}
